\section{The Glory of God}

\begin{quotex}
The essential American soul is hard, stoic, isolate, and a killer. It has never yet melted. 
\flright{\textsc{DH Lawrence}}
\end{quotex}

If it is true that the ``Great Divide\footnote{\url{http://www.gornahoor.net/?p=2244}}'' was the Middle Ages, it might be worthwhile examining to what extent they remain with us, and how we might recover such. Ordinarily, it is claimed that Luther and Calvin and the Reformation ended the Middle Ages, and of course in a purely causal sense (itself limited to certain causes) this is certainly accurate. However, Luther and Calvin are inheritors of the Augustinian tradition, and much of what is praiseworthy in their \emph{traditio} is essentially medieval\footnote{\url{http://www.kinism.net/publications/tkr_summer_2010.pdf}}. Much could be further dug\footnote{\url{http://anglicanrose.wordpress.com/}} out of Reformation history that has gone neglected, yet there is a central strain of Augustinianism which Evola singled out for comment:

\begin{quotex}
As for Christianity in its less popular forms, it presents an aspect of the tragic doctrine of salvation, which to some extent preserves an echo of the ancient truth: the idea--pushed to extremes by Luther and Calvin---that man on earth stands at the crossroads between Salvation and eternal damnation. This point of view, if lived intensely and coherently, could create the conditions for liberation at the moment of death or in post-mortem states.

\flright{\emph{The Hermetic Tradition}, Note 2, page 96}
\end{quotex}


It is not worth arguing that there isn't a lot to object to in Christian theology proper, especially in its generalized concerns and contours, but even in the details. Yet the Augustinian tradition, with both its ``great rewards and great dangers'' (Khuenhelt-Ledhin), was a historical vehicle of impressing upon many men the necessity of pleasing God rather than men (the exoteric phrasing of its content). It gave a great cast of character to certain practicioners which went beyond the mere philosophical fatalism or negation of sensual pleasure which might be predicted by a superficial reading, and this was indeed a result that many lesser adherents naturally devolved. Certainly, Calvin's view (for example) might lead to a mechanical clockwork universe, in which every detail is ``restlessly and ceaselessly ordered'' by the direct hand of God, denying the possibility of actualization of being by the individual, and tending towards a complacent secularism. The anti-intellectual tendencies of Luther are well known.

However, these men were ``medieval'' in their grasp of life as a fairy tale that was dark in the middle, through which men must pass -- hence the emphasis upon transcendence and perseverance, the ``trials of faith''.

Luther himself argued that man has to pass through \emph{Anfechtung}, which was the crushing of the personality under the weight of sin. We can here posit that included in this ``sin'' was the ``ego''. William James investigates this lineage in his \emph{Varieties of Religious Experience}, the tradition of being ``twice-born'' and the necessity of regenerating being.

Calvin's emphasis is different\footnote{\url{http://paulhelmsdeep.blogspot.com/2009/07/augustine-and-calvin-on-knowledge-of.html}}, but his proteges actually proclaimed that in ``order to be saved, one has to be willing to be damned for God's glory''. This paradox is also a crushing of the ego, and an attempt to impress man with the vertical energies which alone can deliver him from purely psychic or all too human phenomena.

Augustine, a mentor to both, situated the locus of this tradition in the ``God-shaped hole'' which exists in the inner heart, and argues that the purely natural heart ``is restless until it rests in Thee''. Although his \emph{Confessions} helped inaugurate the Western odyssey, he remained thoroughly traditional in his attempt to preserve Plato's forms within the mind of God.

As Evola's quote indicates, it was in the intensity of their final focus where the clue for real possibilities lay -- by driving man outside of himself, and forcing the ego to ``melt'' in the presence of real transcendence (a transcendence guaranteed in the magisterial side of the Augustinian tradition, and linked with the ``way of the heart'' as a guide in the dark) the opening up of real religious vision could be made more fully present. Augustinianism in fact gave to the Protestant West whatever freedom from the ego it possessed, and a presence of the Super-Ego, which it could not derive from its geographic situation and cultural history.

\emph{It is here qualified that such would more or less create the possibilities for transformation or pre-transformation over time, including the ``after-life'', rather than fully actualizing such in a truly self-conscious manner. We are not claiming that they represent a coherent perennial tradition -- their coherence is derived from their intensity and whatever echoes they preserve, as well as whatever inspiration they manage to possess.}

In this path is found the meaning of ``original sin'' in the Christian tradition -- an assault upon the accidents and hopelessness of the perverted ego. It was certainly never meant to be applied to the real self, and indeed in moments of transport or inspiration, never was, but almost always only in the ``proclamation of such'' to the ``masses''.

Even if it was muddled, traditionalists are in a position to recognize that mens' experience of an event is one thing, their accurate interpretation and communicating of it something else. This allows them to begin to decode the record of spiritual experiences, even when it seems a wasteland. This is particularly important in Western studies and the reconstitution of our tradition.

Our verdict would be similar to that of Evola's on esoteric Catholicism -- if the Anglo freethinker would return to magisterial Augustinianism, this would not be a regress, but a positive spiritual good. As disrupted as that tradition is on the surface, we have reason to believe that deeper patterns were at work. For instance, Luther drew on the Theologica Germanica\footnote{\url{http://en.wikipedia.org/wiki/Theologia_Germanica}}, which is attributed possibly to Eckhart and associated with the Teutonic knightly order.

\begin{quotex}
Next to the Bible\footnote{\url{http://en.wikipedia.org/wiki/Bible}} and St. Augustine, no book has ever come into my hands from which I have learned more of God and Christ, and man and all things that are.
\end{quotex}


The book advocates the union of the individual with God as an earthly goal. The Divine Drama, in which man is at the center stage, and must choose a path which either damns him to the pits, or leads him to better than Elysium, is given an existential urgency and self-evidential character, along with whatever of ancient doctrine the Western mind was capable of absorbing at that time.

If Gurdjieff's ``way of the heart\footnote{\url{http://en.wikipedia.org/wiki/Fourth_Way}}'' characterizes the Christian path, how much more does it characterize the Augustinian tradition, which is essentially a theological poetry focused on the lone person, nevertheless sustained by God's sole Providence? It may be that in this existential sense, in the ``echo'' (as Evola has it) it might be possible for man to fulfill the Biblical mandate (itself a stronger echo) that each man would ``feel after God, if happily they might find Him'' (Aratus, quoted by Paul). The intense Biblicism and personal piety which characterized the initial Reformation was in this sense an attempt to ``see through the glass darkly'', and to reclaim the lost birthright. In divine sovereignty, all these men saw with Dante that here ``Love and Fate are one''. If they garbled the rapturous report, it is up to others to retune it.

Such investigations into Western history (heretofore interpreted largely by the detritus intelligentsia class, for whom the world before Voltaire was without grace or merit) will lead to surprising avenues of inquiry, in which it becomes clearer and clearer that the rivulet of traditional influence in the West had an odd habit\footnote{\url{http://www.brusselsjournal.com/node/3732}} of showing up, again and again, in strange guises.

The Augustinian tension between the soul that ``believes in order to understand'' and the God Who sovereignly calls such a soul to even desire it echoes the tension in traditional thought between the man who passes through the initiation of his own volition and in full passibility, in order to integrate himself with an order which is eternal and which knows nothing of such vulnerabilities to contingency.

From death, to life. This perennial wisdom may not have ever had a chance to become explicit in the West, but we find it even in our poets:

\begin{quotex}
Take, he said, the belief in immortality, which, according to some men, is a matter of mild indifference. It is really a belief which affects our whole conception of the human race. Consider, he said, the carnage of war, with its pile of unnumbered corpses. It must make some matter to us whether, according to our serious belief, each man has died like a dog, and left nothing in the way of a personal existence behind him, or ``whether out of every Christian-named portion of that ruinous heap there has gone forth into the air and the dead-fallen smoke of battle some astonished condition of soul unwillingly released.''

\flright{\textsc{John Ruskin}, quoted in W. H. Mallock's \textit{Memoirs of Life and Literature}\footnote{\url{http://www.archive.org/details/memoirsoflifelit00mall}}}
\end{quotex}


Is this not precisely Evola's teaching? Is it not time we reclaim all of our birthright?\footnote{Please see \textit{Greeks, Romans, Christians -- Essays in Honor of Abraham Malherbe}, particularly the essay ``The Aeropagus Speech: And Appeal to the Stoic Historian Posidonius against Later Stoics and the Epicureans'', by David Balch. This volume contains excellent explorations of Pauline attitudes towards Stoicism.}


\flrightit{Posted on 2011-04-23 by Logres}

\begin{center}* * *\end{center}

\begin{footnotesize}
\begin{sffamily}
\texttt{James O'Meara on 2011-04-23 at 10:32 said:}

Not to go off on a tangent, but I found the book, or at least the review of it, that you link to quite interesting, since it argues against the view, which I was taught quite explicitly in studying Plotinus, that Greek culture, or at least the works of P, A, and Pl, were preserved by the Arabs and the texts and translations made their way to Europe around the XIIth century.

On the other hand, doesn't the author and review, who start off by mocking the ``Enlightenment'' view of the ``Dark Ages'', lapse into exactly the same Modernist dogma, praising Athens and supposedly Jerusalem for `freedom', `democracy,''curiosity' and all the other motifs Guenon would identify as the roots of the West's deviation? In Evola's terms, they seem to be taking the Church side of the ``Roman Church'' rather than the Roman, but then mis identifying it with the Classical world as such, and contrasting it to Islam. Since Evola preferred the Roman side, he would, and did, prefer Islam, precisely in the ways that the authors deride it!

If, as the authors say, Islam is evil for stoning the Woman Taken in Adultery, while Christ forgives her, is that not exactly what Evola would condemn: an unworldly `spirituality' not fit for constructing a real society. Hence the need to take over the Roman machinery, transforming primitive Christianity into the Roman Church.

Or as Schuon said, Christianity is the ``in a sense'' illegitimate exotericizing of the esoteric core of Judaism, which imbalance Islam is ``providentially'' revealed to compensate and replace? Speaking of Schuon, the authors make much of Islam as ``unable'' to ``assimilate'' the Logos, which somehow prevents it from favoring democracy ! Now, Schuon established this parallel:

Gabriel --- Word spermatic logos --- Mary virgin --- Christ

Gabriel --- Word ``Hear!'' --- Mohammed illiterate --- Koran

I find this a lovely bit of structuralism, and I've often used it to illustrate what G. would call `metaphysical insight', perceiving the hidden identity of the elements of disparate traditions, or what Spengler called `physiognomic tact'' and I resent the authors dirtying it up with their Islamophobia!

\hfill

\texttt{logres on 2011-04-23 at 10:44 said:}

Actually, I should have picked a better site to link to for the book, which was the main point of my link, and which I find fascinating. I would agree (and you are right to point out that more importantly, Schuon and others would agree) that Islam has an important traditionalist role -- that said, their current state of affairs is somewhat different than even a 100 years ago. Wasn't it Guillame Faye (Archaeo-Futurism) who has argued that Europe's choice is between re-imposing real order on themselves, or having it done for them by Islam? So yes, that would be a Providential role. I am not familiar enough with Schuon's writings to comment, except to say that he has an enormous amount to teach.

Even ``conservative'' sites like Brussels Journal have imbibed the Enlightenment mindset (most conservatives simply are not). But it is definitely interesting that Western culture may have maintained far more integrity in various ways than we give it credit for. It's possible that the Celts had preserved some of the texts as well -- have to look into Erigena.

\hfill

\end{sffamily}
\end{footnotesize}

